\documentclass{article}
\usepackage[UTF8]{ctex}
\usepackage{amsmath,amssymb,amsfonts}
\usepackage{enumitem}
\usepackage{geometry}
\geometry{a4paper, margin=1in}
% === 错题本排版优化（不参与题目标识） ===
\setlength{\parindent}{0pt}
\setlength{\parskip}{0.35em plus 0.10em minus 0.08em}
\setlength{\emergencystretch}{3em}
\setlist{itemsep=0.25em, topsep=0.40em, parsep=0pt}
\allowdisplaybreaks[2]
% === 错题本排版优化结束 ===
\begin{document}
\section*{四川大学2024年数学分析真题}
\section*{165．四川大学2024年数学分析试卷（回忆版）}

\subsection{题目1}
1． \(\lim _{n \rightarrow \infty} \sqrt[n]{1+\sqrt[n]{2+\cdots+\sqrt[n]{n}}}\) ．

\subsection{题目2}
2． \(\lim _{x \rightarrow 0}\left(\frac{x}{\sin x}-\cos x+1\right) \frac{1}{x}\) ．

\subsection{题目3}
3． \(\lim _{n \rightarrow \infty} \frac{\ln 2+\ln 3+\cdots+\ln n}{n^{2} \ln n}\) ．

\subsection{题目4}
4． \(\lim _{x \rightarrow 0^{+}} \frac{x}{\cos x} \int_{x^{-1}}^{x^{-2}}\left(1+\frac{1}{2 t}\right)^{t} \sin \frac{1}{\sqrt{t}} \mathrm{~d} t\) ．

\subsection{题目5}
1． \(\int \frac{1+x}{1-x^{4}} \mathrm{~d} x\) ．

\subsection{题目6}
2．\(\oint_{L} \frac{y}{(2-x)^{2}+y^{2}} \mathrm{~d} y+\frac{x}{(2-x)^{2}+y^{2}} \mathrm{~d} y\) ，其中 \(L\) 是不经过 \((2,0)\) 点的封闭曲线，方向取逆时针．

\subsection{题目7}
3． \(\iint_{\sum} \frac{2}{x} \mathrm{~d} y \mathrm{~d} z+\mathrm{d} z \mathrm{~d} x+\frac{1}{z} \mathrm{~d} x \mathrm{~d} y\) ，其中 \(\sum\) 是 \(x^{2}+y^{2}+z^{2}=1\) 的表面，方向取外侧．

\subsection{题目8}
4． \(\int_{0}^{+\infty} \frac{e^{-a x}-e^{-b x}}{x} \mathrm{~d} x\) ．

三、证明区间上有第一类间断点的函数一定没有原函数．

四．设 \(f(x)\) 在 \((0,1)\) 上有界且可导， \(\lim _{x \rightarrow 0^{+}} f(x)\) 不存在，证明一定存在数列 \(\left\{x_{n}\right\}\) ，当 \(\lim _{n \rightarrow \infty} x_{n}=0\) 时一定有 \(\lim _{n \rightarrow \infty} f\left(x_{n}\right)=0\)

五．设 \(F(x, y)=0\) 是平面上一条八字型光滑曲线，说明曲线上满足 \(\frac{\partial F}{\partial x}=\frac{\partial F}{\partial y}=0\) 的点至少有几个兵说明理由．

六．连续可微函数 \(z=z(x, y)\) 由 \(F(x z-y, x-y z)=0\) 唯一确定．

\begin{itemize}
  \item[（1）] 求 \(\frac{\partial z}{\partial x}, \frac{\partial z}{\partial y}\) 。
\end{itemize}

\subsection{题目6}
原题边界“（2）”：证明：\(\oint_{L}\left(x y^{2}+2 y z\right) \mathrm{d} y-\left(2 x z+y z^{2}\right) \mathrm{d} x=2 \pi, L\) 为 \(x^{2}+y^{2}=1\) ，方向取逆时针．

七．讨论反常积分 \(\int_{0}^{+\infty} \frac{\ln x}{x} \sin x \mathrm{~d} x\) 的收敛性（包括条件收敛和绝对收敛）．

八．证明 \(f(x, y, z)=\frac{1}{1-x y z}\) 在 \([0,1) \times[0,1) \times[0,1)\) 上非一致连续．

九．设 \(x<f_{1}(x)<\sqrt{x}, f_{n}(x)=\sqrt{f_{n-1}(x)}\) ．

\begin{itemize}
  \item[（1）] 对任意的 \(x_{0} \in(0,1)\) ，证明：\(\left\{f_{n}\left(x_{0}\right)\right\}\) 是单调有界数列．
  \item[（2）] 证明：\(\left\{f_{n}(x)\right\}\) 在 \((0,1)\) 上一致收敛．
\end{itemize}

十．想不起来了．
\end{document}
